\documentclass[11pt]{article}
\usepackage{metricsstyle}
\usepackage{array,booktabs}
\usepackage[most]{tcolorbox}
\definecolor{qnavy}{RGB}{31,58,96}
\definecolor{whygreen}{RGB}{233,244,236}
\definecolor{whyline}{RGB}{60,130,80}
\newtcolorbox{why}{enhanced, breakable, colback=whygreen, colframe=whygreen,
  borderline west={2pt}{0pt}{whyline}, boxrule=0pt, arc=0pt,
  left=6pt, right=4pt, top=3pt, bottom=3pt, before skip=4pt, after skip=6pt}
\newcommand{\Aa}{R(X'X)^{-1}R'}
\tikzset{fb/.style={draw=qnavy!70,fill=qnavy!5,rounded corners=2pt,align=center,
  inner sep=4pt,text width=4.6cm,font=\small},
  fbw/.style={draw=qnavy!70,fill=qnavy!5,rounded corners=2pt,align=center,
  inner sep=4pt,text width=11cm,font=\small},
  fbg/.style={draw=whyline,fill=whygreen,rounded corners=2pt,align=center,
  inner sep=4pt,text width=11cm,font=\small},
  ar/.style={-{Stealth[length=4pt]},line width=0.6pt}}

\begin{document}
\begin{center}
  {\large\scshape Econometrics I --- Midterm Review}\\[3pt]
  {\LARGE\bfseries The $F$ Statistic: One Picture}\\[4pt]
  {\itshape How $R\beta=r$ leads to every form of $F$ in Lessons 5--7}
\end{center}
\vspace{0.2em}\hrule\vspace{0.8em}

\section{The picture}
All the $F$ formulas are \emph{one} statistic. It measures a single quantity, the gap
$Rb-r$ between what the data say and what $H_0$ claims. There are three ways to
compute the same numerator, depending on which model you have estimated.

\begin{center}
\begin{tikzpicture}
  \node[fbw] (A) at (0,0) {\textbf{$H_0:\ R\beta=r$}\quad ($R$ is $q\times K$: $q$ restrictions)};
  \node[fbw] (B) at (0,-1.5) {\textbf{The gap:} $Rb-r=R(X'X)^{-1}X'u$ under $H_0$\\[2pt]
     $\Var(Rb-r)=\sigma^{2}A,\qquad A=R(X'X)^{-1}R'$ \hfill (L5)};
  \node[fb] (W) at (-5.2,-3.9) {\textbf{Wald} (fit unrestricted only)\\[2pt]
     $(Rb-r)'A^{-1}(Rb-r)$\\ \footnotesize Lesson 5};
  \node[fb] (S) at (0,-3.9) {\textbf{Lagrange multiplier} (fit restricted only)\\[2pt]
     $\lambda_*'A\,\lambda_*$\\ \footnotesize Lesson 7};
  \node[fb] (R) at (5.2,-3.9) {\textbf{SSR comparison} (fit both)\\[2pt]
     $e_*'e_*-e'e=RSSR-USSR$\\ \footnotesize Lesson 6};
  \node at (-2.6,-3.9) {$=$}; \node at (2.6,-3.9) {$=$};
  \node[fb] (C1) at (-2.9,-6.2) {numerator$/\sigma^{2}\ \sim\ \chi^{2}_{q}$};
  \node[fb] (C2) at (2.9,-6.2) {$e'e/\sigma^{2}=u'Mu/\sigma^{2}\ \sim\ \chi^{2}_{n-K}$};
  \node[fbg] (F) at (0,-7.9) {$\displaystyle F=\frac{\text{numerator}/q}{e'e/(n-K)}\ \sim\ F(q,\,n-K)$\\[3pt]
     \footnotesize $\sigma^{2}$ cancels; the two $\chi^{2}$ are independent because $MN=0$ (L5)};
  \draw[ar] (A) -- (B);
  \draw[ar] (B.south) -- ++(0,-0.35) -| (W.north);
  \draw[ar] (B.south) -- ++(0,-0.35) -| (R.north);
  \draw[ar] (B.south) -- (S.north);
  \draw[ar] (W.south) -- ++(0,-0.4) -| (C1.north);
  \draw[ar] (S.south) -- (C1.north east);
  \draw[ar] (R.south) -- ++(0,-0.4) -| (C1.north);
  \draw[ar] (C1.south) |- ++(0,-0.45) -| (F.north);
  \draw[ar] (C2.south) -- ++(0,-0.45) -| (F.north);
\end{tikzpicture}
\end{center}

\begin{why}
\textbf{Why it works:} a large gap means $H_0$ forces $b$ far from where the data put it.
Forcing it costs fit (RSSR $\gg$ USSR), makes the multiplier large, and makes the Wald
distance large. These are three descriptions of the same event.
\end{why}

\section{Step by step, as in the notes}

\topic{Step 1 (Lesson 5): the gap and its distribution}
Under $H_0$,
\[ Rb-r=R\beta+R(X'X)^{-1}X'u-r=R(X'X)^{-1}X'u,\qquad
   \Var(Rb-r)=\sigma^{2}R(X'X)^{-1}R'. \]
Measure the distance by the quadratic form weighted by the inverse covariance:
\[ \frac{1}{\sigma^{2}}(Rb-r)'\bigl[R(X'X)^{-1}R'\bigr]^{-1}(Rb-r)=\frac{1}{\sigma^{2}}u'Nu\sim\chi^{2}_{q},\]
with $N=X(X'X)^{-1}R'\bigl(R(X'X)^{-1}R'\bigr)^{-1}R(X'X)^{-1}X'$ symmetric and idempotent,
$\operatorname{rank}N=\operatorname{tr}N=q$.

\topic{Step 2 (Lesson 5): the other $\chi^{2}$ and independence}
\[ \frac{e'e}{\sigma^{2}}=\frac{u'Mu}{\sigma^{2}}\sim\chi^{2}_{n-K},\qquad MN=0\ \text{(since }MX=0)\ \Rightarrow\ \text{independent.}\]

\topic{Step 3 (Lesson 5): form the ratio}
\[ F=\frac{\frac{1}{\sigma^{2}}u'Nu/q}{\frac{1}{\sigma^{2}}u'Mu/(n-K)}
    =\frac{(Rb-r)'\bigl[R(X'X)^{-1}R'\bigr]^{-1}(Rb-r)/q}{s^{2}}\sim F(q,n-K). \]
This is the \textbf{master formula}. All later forms are rewrites of its numerator.
For $q=1$ it is $t^{2}_{n-K}$ ($F(1,q)=t_q^{2}$ in the notes).

\topic{Step 4 (Lesson 6): the numerator as a difference in SSR}
Minimise $e_*'e_*$ subject to $Rb_*=r$ (Lagrangian). The first-order conditions give
\[ \lambda_*=\bigl[R(X'X)^{-1}R'\bigr]^{-1}(r-Rb),\qquad
   b_*=b-(X'X)^{-1}R'\bigl[R(X'X)^{-1}R'\bigr]^{-1}(Rb-r). \]
The restricted residual is $e_*=e-X(b_*-b)$, and $X'e=0$ kills the cross term:
\[ e_*'e_*=e'e+(b_*-b)'X'X(b_*-b). \]
Substituting $b_*-b$, the $X'X$ cancels against the two $(X'X)^{-1}$:
\[ \boxed{\,e_*'e_*-e'e=(Rb-r)'\bigl[R(X'X)^{-1}R'\bigr]^{-1}(Rb-r)\,} \]
Hence
\[ F=\frac{(e_*'e_*-e'e)/q}{e'e/(n-K)}=\frac{(RSSR-USSR)/q}{USSR/(n-K)}. \]

\topic{Step 5 (Lesson 6): subset $H_0:\beta_1=0$ via FWL}
With $R=[I_{k_1}\ \ 0]$, $q=k_1$ and FWL, $b_1\sim N(0,\sigma^{2}(X_1'M_2X_1)^{-1})$, so
\[ F=\frac{b_1'X_1'M_2X_1b_1/k_1}{e'e/(n-K)}. \]
Premultiply $y=X_1b_1+X_2b_2+My$ by $M_2$ (kills $X_2$, $M_2M=M$):
$M_2y=M_2X_1b_1+My$, so $y'M_2y=b_1'X_1'M_2X_1b_1+y'My$, and
\[ b_1'X_1'M_2X_1b_1=y'M_2y-y'My=RSSR-USSR. \]
Here $M_2y$ is the residual of the regression with $H_0$ imposed (only $X_2$).

\topic{Step 6 (Lesson 7): the LM view}
Under $H_0$, $\lambda_*$ has mean $0$ and $\Var(\lambda_*)=\sigma^{2}\bigl[R(X'X)^{-1}R'\bigr]^{-1}$, so
\[ \lambda_*'\bigl[\Var(\lambda_*)\bigr]^{-1}\lambda_*=\frac{1}{\sigma^{2}}(r-Rb)'\bigl[R(X'X)^{-1}R'\bigr]^{-1}(r-Rb). \]
This is identical to the Wald quantity, since $\lambda_*$ is $Rb-r$ times a fixed nonsingular matrix.
Testing the multiplier (restricted fit) is the same test as testing the gap (unrestricted fit).

\section{The family of formulas}
\small
\renewcommand{\arraystretch}{1.5}
\begin{tabular}{@{}>{\raggedright\arraybackslash}p{2.9cm}>{\raggedright\arraybackslash}p{7.6cm}>{\raggedright\arraybackslash}p{4.6cm}@{}}
\toprule
\textbf{Name} & \textbf{Formula} & \textbf{Use when} \\
\midrule
Master (L5) & $\dfrac{(Rb-r)'[R(X'X)^{-1}R']^{-1}(Rb-r)/q}{s^{2}}$ & the question gives $b$, $(X'X)^{-1}$, $R$, $r$ \\
SSR form (L6) & $\dfrac{(RSSR-USSR)/q}{USSR/(n-K)}$ & you have (or can run) both regressions \\
Subset, FWL (L6) & $\dfrac{b_1'X_1'M_2X_1b_1/k_1}{e'e/(n-K)}=\dfrac{(y'M_2y-y'My)/k_1}{y'My/(n-K)}$ & $H_0:\beta_1=0$, proofs \\
LM (L7) & $\dfrac{\lambda_*'[\Var(\lambda_*)]^{-1}\lambda_*/q}{s^{2}}$ with $\sigma^{2}$ replaced & restricted fit only is available \\
$t^{2}$ (L5) & $F(1,n-K)=t_{n-K}^{2}$ & $q=1$ \\
Wald, $\sigma^{2}$ known & $\dfrac{RSSR-USSR}{\sigma^{2}}\sim\chi^{2}_q$ & $\sigma^{2}$ given; no denominator \\
\midrule
\multicolumn{3}{@{}l}{\textit{Special cases (derived from the SSR form; not separate theorems)}}\\
Overall significance & $\dfrac{ESS/(k-1)}{RSS/(n-k)}=\dfrac{R^{2}/(k-1)}{(1-R^{2})/(n-k)}$ & all slopes $=0$: restricted model is constant only, $RSSR=TSS$ \\
$R^{2}$ form & $\dfrac{(R_U^{2}-R_R^{2})/q}{(1-R_U^{2})/(n-K)}$ & same $y$ in both models (same $TSS$) \\
Chow & $\dfrac{(SSR_p-SSR_1-SSR_2)/k}{(SSR_1+SSR_2)/(n_1+n_2-2k)}$ & stability; $R=[I\ \ {-I}]$, restricted $=$ pooled \\
\bottomrule
\end{tabular}
\normalsize

\begin{why}
\textbf{ESS here:} for overall significance, ESS means $y'(P_X-P_\iota)y=TSS-RSS$ (centred, with a
constant). The notes' uncentred $y'P_Xy$ would be the wrong numerator. ESS appears only in that
one special case; for a general restriction use $RSSR-USSR$.
\end{why}

\section{Why the denominator is the \emph{unrestricted} $s^{2}$ (Lesson 7)}
Using $RSSR/(n-K)$ in the denominator breaks independence. For the subset case,
\[ M(M_2-M)=M-M=0\ \Rightarrow\ \text{independent (an $F$)},\qquad
   M_2(M_2-M)=M_2-M\ne0\ \Rightarrow\ \text{dependent}, \]
and $u'(M_2-M)u/u'M_2u$ has the form $A/(A+B)$, i.e.\ a $\mathrm{Beta}\bigl(\tfrac{k_1}{2},\tfrac{n-K}{2}\bigr)$,
not an $F$.

\section{Exam recipe}
\begin{enumerate}[label=\arabic*.]
  \item Write $H_0$ as $R\beta=r$ and count $q$ (rows of $R$).
  \item Given $b$, $(X'X)^{-1}$: use the \textbf{master} form. Given sums of squares: use \textbf{RSSR/USSR}.
  \item USSR is from the model \emph{without} $H_0$; RSSR is from the model with $H_0$ \emph{substituted in}.
  \item $df_2=n-K$ of the unrestricted model. Compare to $F_{\alpha}(q,n-K)$ and state reject / fail to reject.
\end{enumerate}
\end{document}
